\begin{abstract}

L'efficacia e la scalabilità di tre algoritmi di Graph Machine Learning vengono analizzate applicandoli a un knowledge graph biomedico di grandi dimensioni. È stato scelto il dataset PharMeBiNet V3 \cite{konigs2022pharmebinet}, una rete eterogenea con circa 8 milioni di nodi distribuiti su oltre 80 tipi e 53 milioni di relazioni distribuite su 410 tipi distinti, che interconnette composti chimici, geni, proteine e patologie provenienti da oltre 40 banche dati farmacologiche e biologiche.

Il problema studiato è la link prediction tra nodi \textit{Chemical} e \textit{Disease}, biologicamente rilevante nel contesto del drug repurposing, ovvero prevedere relazioni terapeutiche non ancora documentate tra farmaci esistenti e malattie note. Il sottografo estratto per gli esperimenti comprende 594.343 nodi Chemical, 23.580 nodi Disease e 123.361 relazioni dirette tra loro. I tre algoritmi confrontati sono GraphSAGE, TransE e HashGNN, scelti per le loro diverse strategie di rappresentazione, rispettivamente: omogenea con aggregazione dei vicini, traslazionale nello spazio vettoriale, ed eterogenea con proiezioni differenziate per tipo di nodo.

L'infrastruttura sperimentale è un cluster Neo4j Enterprise 2026.04.0 distribuito su quattro macchine virtuali cloud, nella configurazione di 1 primary e 3 secondary. Il benchmark di scalabilità valuta l'impatto del numero di nodi attivi sui tempi di estrazione dei dati verso la pipeline PyTorch Geometric, variando la configurazione tra 1, 2 e 4 nodi.

I risultati mostrano che GraphSAGE raggiunge la precisione più alta con un AUC-ROC medio di 0.993 su cinque run indipendenti, mentre TransE si rivela il modello più efficiente, registrando tempi di training circa tre volte inferiori pur mantenendo una qualità predittiva eccellente (AUC-ROC > 0.98). Sul fronte infrastrutturale, il cluster distribuito non produce speedup significativi su query sequenziali di estrazione batch, evidenziando come il collo di bottiglia sia legato all'overhead di coordinamento di rete piuttosto che alla disponibilità di risorse computazionali.
\end{abstract}